\documentclass[12pt]{article}
\usepackage[utf8]{inputenc}
\usepackage[T1]{fontenc}
\usepackage[margin=1in]{geometry}
\usepackage{setspace}
\usepackage{xcolor}
\usepackage[colorlinks=true, urlcolor=blue!60!black, linkcolor=blue!60!black]{hyperref}
\setlength{\parindent}{2em}
\setlength{\parskip}{1.25em}
\renewcommand{\baselinestretch}{1.0}
\begin{document}
\begin{center}
{\LARGE \textit{Board of Trustees of the University of Alabama v. Garrett}, 531 U.S. 356 (2001)\par}
\vspace{0.8em}
{\large Case Brief\par}
\end{center}
\vspace{1em}

\section{Memory Jogger}

In \href{https://www.courtlistener.com/opinion/118407/board-of-trustees-of-univ-of-ala-v-garrett/}{\textit{Board of Trustees of the University of Alabama v. Garrett}}, the Supreme Court held 5--4 that the Eleventh Amendment bars state employees from suing their state employers in federal court for money damages under Title I of the Americans with Disabilities Act (ADA), because Congress failed to document a sufficient pattern of unconstitutional state employment discrimination and enacted remedies not ``congruent and proportional'' under Section 5 of the Fourteenth Amendment.

\section{Detailed Case Facts}

This consolidated litigation involved two Alabama state employees who alleged disability discrimination in employment under Title I of the Americans with Disabilities Act of 1990 (ADA), 42 U.S.C. \S\S~12111--12117 (\href{https://supreme.justia.com/cases/federal/us/531/356/}{531 U.S. at 360--362}).

First, respondent Patricia Garrett, a registered nurse, served as Director of Nursing, OB/Gyn/Neonatal Services, at the University of Alabama in Birmingham Hospital. Diagnosed with breast cancer in 1994, Garrett underwent a lumpectomy, radiation therapy, and chemotherapy, requiring substantial medical leave. When she returned to work in July 1995, her supervisor told her she would have to give up her Director position, and Garrett transferred to a lower-paying nurse-manager role.

Second, respondent Milton Ash worked as a security officer for the Alabama Department of Youth Services. Upon starting the job, Ash informed the Department that he suffered from chronic asthma and, on his physician's advice, asked to have his duties modified to minimize exposure to carbon monoxide and cigarette smoke. Later diagnosed with sleep apnea, Ash also requested reassignment to daytime shifts per his doctor's recommendation. The Department denied all requested accommodations, and shortly after Ash filed an Equal Employment Opportunity Commission (EEOC) discrimination charge, his performance evaluations dropped below his prior ratings (\href{https://tile.loc.gov/storage-services/service/ll/usrep/usrep531/usrep531356/usrep531356.pdf}{531 U.S. at 362}).

Title I of the ADA prohibits covered employers---explicitly including States---from discriminating against a ``qualified individual with a disability'' in hiring, advancement, discharge, compensation, and other terms of employment (42 U.S.C. \S~12112(a)). It requires employers to make ``reasonable accommodations'' (such as modifying facilities, schedules, or equipment) unless doing so imposes an ``undue hardship'' (\S~12112(b)(5)(A); \S~12111(9)), and forbids standards or methods of administration that have a disparate impact on disabled individuals (\S~12112(b)(3)(A)). Congress expressly sought to abrogate state sovereign immunity in 42 U.S.C. \S~12202 and invoked both its Commerce Clause power and its enforcement authority under Section 5 of the Fourteenth Amendment (\S~12101(b)(4)).

Historically and politically, the ADA was enacted in 1990 with overwhelming bipartisan support and signed by President George H. W. Bush (who filed a supportive amicus statement in \textit{Garrett}) after 13 congressional hearings and a nationwide Task Force on the Rights and Empowerment of Americans with Disabilities. By 2001, however, the Rehnquist Court's ``federalism revolution'' had sharply curtailed Congress's power to subject nonconsenting States to private damages suits: \textit{Seminole Tribe of Florida v. Florida}, 517 U.S. 44 (1996), barred abrogation under Article I powers; \textit{City of Boerne v. Flores}, 521 U.S. 507 (1997), required Section 5 legislation to exhibit ``congruence and proportionality'' to documented Fourteenth Amendment violations; and \textit{Kimel v. Florida Board of Regents}, 528 U.S. 62 (2000), invalidated abrogation under the Age Discrimination in Employment Act. \textit{Garrett} tested whether the ADA's far more extensive legislative record could survive that heightened Section 5 scrutiny.

\section{Procedural History}

\begin{itemize}
\item \textbf{U.S. District Court for the Northern District of Alabama:} Garrett and Ash filed separate federal lawsuits seeking money damages under the ADA against their respective Alabama state employers. In a single opinion disposing of both actions, \textit{Garrett v. Board of Trustees of the University of Alabama}, 989 F. Supp. 1409 (N.D. Ala. 1998), the District Court granted summary judgment for the State, holding that Congress exceeded its constitutional authority in attempting to abrogate Eleventh Amendment immunity under the ADA (\href{https://supreme.justia.com/cases/federal/us/531/356/}{531 U.S. at 362--363}).
\item \textbf{U.S. Court of Appeals for the Eleventh Circuit:} Consolidating the appeals, the Eleventh Circuit reversed in \textit{Garrett v. University of Alabama at Birmingham Board of Trustees}, 193 F.3d 1214 (11th Cir. 1999), adhering to circuit precedent holding that the ADA validly abrogated state sovereign immunity.
\item \textbf{U.S. Supreme Court:} The Court granted certiorari, 529 U.S. 1065 (2000), to resolve a circuit split over private damages suits against States under the ADA. Oral argument took place on October 11, 2000, and the Court issued its decision on February 21, 2001. By a 5--4 vote, the Court \textbf{reversed} as to Title I and dismissed as improvidently granted the portion of the writ concerning Title II (public services), because the parties had not briefed whether Title II covers employment claims (\href{https://www.law.cornell.edu/supremecourt/text/531/356}{531 U.S. at 360 n. 1, 374}).
\end{itemize}

\section{Judicial Votes}

The Court divided 5--4 along its familiar Rehnquist-era federalism alignment:

\begin{itemize}
\item \textbf{Majority (5):} Chief Justice Rehnquist (author), joined by Justices O'Connor, Scalia, Kennedy, and Thomas.
\item \textbf{Concurrence (2):} Justice Kennedy (author), joined by Justice O'Connor.
\item \textbf{Dissent (4):} Justice Breyer (author), joined by Justices Stevens, Souter, and Ginsburg.
\end{itemize}

\section{Holding}

\begin{itemize}
\item \textbf{Judgment:} Reversed. Private suits in federal court by state employees seeking money damages against a State for failure to comply with Title I of the ADA are barred by the Eleventh Amendment.
\item \textbf{Rule Announced:} Congress may abrogate state sovereign immunity only when it unequivocally expresses its intent to do so and acts pursuant to a valid exercise of Section 5 of the Fourteenth Amendment (not Article I). To qualify as valid Section 5 legislation enforcing a right subject only to rational-basis review (such as disability classifications under \textit{City of Cleburne v. Cleburne Living Center, Inc.}, 473 U.S. 432 (1985)), Congress must (1) identify a history and pattern of unconstitutional discrimination by the \textit{States themselves} (excluding private actors and local governments), and (2) enact remedies that are ``congruent and proportional'' to the targeted constitutional violations (\href{https://supreme.justia.com/cases/federal/us/531/356/}{531 U.S. at 363--365, 368--374}).
\item \textbf{Holding vs. Dicta:} The binding holding is confined to private federal-court suits for \textit{money damages} against States under \textit{Title I} of the ADA. In footnote 1, the Court expressly declined to decide whether Title II validly abrogates immunity (later upheld for courtroom access in \href{https://supreme.justia.com/cases/federal/us/541/509/}{\textit{Tennessee v. Lane}, 541 U.S. 509 (2004)}). In footnote 9, Chief Justice Rehnquist noted in influential dicta that Title I's substantive standards still bind the States and remain enforceable by the United States in damages suits, by private individuals seeking prospective injunctive relief under \textit{Ex parte Young}, 209 U.S. 123 (1908), and through state disability-rights laws.
\end{itemize}

\section{Analysis of Opinions}

\subsection{Majority Opinion (Chief Justice Rehnquist)}

Chief Justice Rehnquist begins by restating the Court's sovereign-immunity baseline: although the Eleventh Amendment's text refers to suits by out-of-state or foreign citizens, \textit{Hans v. Louisiana}, 134 U.S. 1 (1890), and \textit{Seminole Tribe} establish that nonconsenting States cannot be sued in federal court by their own citizens (\href{https://supreme.justia.com/cases/federal/us/531/356/}{531 U.S. at 363}). Because Congress's intent to abrogate in 42 U.S.C. \S~12202 was indisputable, and Article I powers cannot abrogate immunity, the sole question was whether Title I was ``appropriate legislation'' under Section 5 of the Fourteenth Amendment (\textit{Fitzpatrick v. Bitzer}, 427 U.S. 445 (1976); \textit{City of Boerne v. Flores}, 521 U.S. 507 (1997)).

Rehnquist structures the Section 5 analysis in three steps:
\begin{enumerate}
\item \textbf{Scope of the Underlying Constitutional Right:} Relying on \textit{Cleburne}, Rehnquist holds that disability is not a suspect or quasi-suspect classification; state action affecting disabled persons triggers only minimum rational-basis review (\href{https://supreme.justia.com/cases/federal/us/531/356/}{531 U.S. at 365--368}). Under rational-basis scrutiny, the challenger must negate ``any reasonably conceivable state of facts'' supporting the classification (\textit{Heller v. Doe}, 509 U.S. 312 (1993)). Although negative attitudes or fear often accompany irrational bias, ``their presence alone does not a constitutional violation make.'' Thus, the Fourteenth Amendment permits States to ``quite hardheadedly---and perhaps hardheartedly---hold to job-qualification requirements which do not make allowance for the disabled.''
\item \textbf{Legislative Record of State Discrimination:} Rehnquist next examines whether Congress documented a history and pattern of \textit{unconstitutional} employment discrimination by \textit{States}. He excludes evidence involving cities and counties because local governments lack Eleventh Amendment immunity (\textit{Lincoln County v. Luning}, 133 U.S. 529 (1890)) (\href{https://supreme.justia.com/cases/federal/us/531/356/}{531 U.S. at 368--369}). Turning to the record, he discounts respondents' ``half a dozen examples'' of state employment incidents and dismisses the roughly 300 state-level anecdotes in Justice Breyer's Appendix C as unexamined submissions to a Task Force---most concerning public services (Title II) rather than employment, and none proved irrational under \textit{Cleburne}. Given that States employed over 4.5 million workers in 1990 and that Senate and House committee reports specifically emphasized discrimination in ``employment in the private sector,'' Rehnquist concludes Congress found no pattern of unconstitutional state employment discrimination.
\item \textbf{Congruence and Proportionality:} Even assuming a pattern existed, Title I's remedies far outstrip Section 1's guarantees (\href{https://supreme.justia.com/cases/federal/us/531/356/}{531 U.S. at 372--374}). Whereas conserving scarce public funds by hiring workers who can use existing facilities is entirely rational under the Equal Protection Clause, the ADA mandates ``reasonable accommodation'' and shifts the burden onto the State to prove ``undue hardship.'' Likewise, Title I bars facially neutral practices with disparate impacts, even though \textit{Washington v. Davis}, 426 U.S. 229 (1976), holds that disparate impact alone does not violate the Equal Protection Clause. Contrasting the ADA with the Voting Rights Act of 1965 upheld in \textit{South Carolina v. Katzenbach}, 383 U.S. 301 (1966), Rehnquist concludes that upholding Title I abrogation would allow Congress to ``rewrite'' \textit{Cleburne}.
\end{enumerate}

\subsection{Concurring Opinion (Justice Kennedy)}

Justice Kennedy, joined by Justice O'Connor, joins the Court's opinion in full while striking a noticeably more sympathetic tone toward the ADA's moral mission (\href{https://supreme.justia.com/cases/federal/us/531/356/}{531 U.S. at 374--376}). Kennedy acknowledges that prejudice against disabled individuals often arises not from malice but from insensitivity, indifference, or instinctive discomfort, and he praises the ADA as a ``milestone on the path to a more decent, tolerant, progressive society'' because ``the law can be a teacher.''

Yet Kennedy draws a sharp constitutional line between societal prejudice and official state constitutional wrongdoing. Charging a State in its sovereign capacity with a pattern of purposeful Equal Protection violations is a grave accusation, particularly when premised on legislative omission rather than hostile action. If States had systematically violated the Fourteenth Amendment, Kennedy argues, ``confirming judicial documentation'' in state and federal court decisions would exist---yet it does not. Stressing that the narrow issue is only whether private individuals may extract money damages from state treasuries without state consent (as opposed to federal enforcement), Kennedy concludes the necessary predicate was not established.

\subsection{Dissenting Opinion (Justice Breyer)}

Justice Breyer, joined by Justices Stevens, Souter, and Ginsburg, delivers a sweeping dissent accompanied by three appendices cataloging 13 congressional hearings, prior federal disability statutes, and roughly 300 state-level discrimination submissions to the congressional Task Force (\href{https://supreme.justia.com/cases/federal/us/531/356/}{531 U.S. at 376--424}).

Breyer attacks the majority on three principal fronts:
\begin{enumerate}
\item \textbf{Treating Congress Like an Administrative Agency or Lower Court:} Breyer argues that Congress assembled a vast record of ``massive, society-wide discrimination,'' including hundreds of state-government incidents, and reasonably inferred that state employers were not immune from pervasive stereotypic assumptions (\href{https://supreme.justia.com/cases/federal/us/531/356/}{531 U.S. at 377--382}). Excluding local governments makes little sense because they are Fourteenth Amendment state actors closely linked to state supervision. More fundamentally, ``a legislature is not a court of law'': Congress need not adjudicate each anecdote to prove the absence of any conceivable rational basis.
\item \textbf{Inverting Judicial Restraint:} Under \textit{Cleburne} and \textit{FCC v. Beach Communications, Inc.}, 508 U.S. 307 (1993), rational-basis review and its presumption of constitutionality are rules of \textit{judicial} restraint reflecting courts' institutional limits---not limits on Congress's Section 5 authority (\href{https://supreme.justia.com/cases/federal/us/531/356/}{531 U.S. at 382--385}). Indeed, \textit{Cleburne} itself expressly deferred to legislative competence in addressing disability discrimination. Applying a judicial-restraint burden to strike down an act of Congress stands the principle ``on its head.''
\item \textbf{Broad Section 5 Remedial Authority and Federalism:} Invoking \textit{Ex parte Virginia}, 100 U.S. 339 (1880), \textit{Katzenbach v. Morgan}, 384 U.S. 641 (1966), and \textit{McCulloch v. Maryland}, 4 Wheat. 316 (1819), Breyer contends that Section 5 empowers Congress to use any rational means to enforce Equal Protection, including reasonable-accommodation and disparate-impact rules (\href{https://supreme.justia.com/cases/federal/us/531/356/}{531 U.S. at 385--389}). He finds it deeply anomalous that the Court applies minimal rational-basis review to laws that \textit{burden} disabled persons while applying near-strict scrutiny to a federal statute designed to \textit{protect} them. Finally, the Civil War Amendments were specifically designed to override state sovereignty; blocking decentralized private damages suits merely pushes Congress toward more intrusive federal agency enforcement and injunctions.
\end{enumerate}

\section{Examples: Future Applications}

\begin{enumerate}
\item \textbf{Same Side (Abrogation Invalid / Sovereign Immunity Bars Suit):} Congress enacts a federal statute prohibiting employment discrimination based on genetic predisposition to non-imminent health conditions, mandating paid preventive-leave accommodations, and authorizing private money-damages suits against States. The legislative record focuses almost exclusively on private employers and insurers, citing only four unverified anecdotes involving state agencies. Under \textit{Garrett}, because genetic status receives only rational-basis review under the Equal Protection Clause and Congress documented neither a pattern of unconstitutional state conduct nor congruent-and-proportional remedies, the Eleventh Amendment bars private damages suits against States.
\item \textbf{Same Side (Abrogation Invalid / Sovereign Immunity Bars Suit):} Congress amends a federal statute to let state employees sue States for money damages whenever a state agency declines to provide paid caregiver leave for an employee's own serious medical condition (the self-care provision, analogous to \textit{Coleman v. Court of Appeals of Maryland}, 566 U.S. 30 (2012)), without documenting that States systematically denied self-care leave on a sex-stereotyped or otherwise unconstitutional basis. Under \textit{Garrett}, absent a targeted record of unconstitutional state discrimination, Congress cannot abrogate Eleventh Amendment immunity.
\item \textbf{Opposite Side (Abrogation Valid):} Congress enacts a federal statute authorizing private money-damages suits against States under Title II of the ADA specifically to enforce disabled citizens' fundamental due-process right of physical access to state courthouses, supported by an extensive legislative record of inaccessible state judicial facilities. As the Court later held in \href{https://supreme.justia.com/cases/federal/us/541/509/}{\textit{Tennessee v. Lane}, 541 U.S. 509 (2004)}, this survives \textit{Garrett} because it enforces a fundamental constitutional right triggering heightened scrutiny rather than mere rational-basis review, and the remedy is narrowly tailored to a documented pattern of state barriers.
\item \textbf{Opposite Side (Abrogation Valid):} Congress enacts a family-care leave mandate authorizing private damages suits against States after compiling a substantial legislative record showing that state employers systematically granted maternity leave to women while denying parallel leave to men based on unconstitutional gender stereotypes. Because sex-based classifications trigger heightened intermediate scrutiny under the Equal Protection Clause, Congress faces a lower evidentiary threshold than in \textit{Garrett} and can validly abrogate state immunity (as in \textit{Nevada Department of Human Resources v. Hibbs}, 538 U.S. 721 (2003)).
\item \textbf{Fence-Sitter (Borderline Case):} Congress passes a statute authorizing private damages suits against States for failing to provide accessible state-university websites and online course portals to blind students, relying on 60 documented state-university complaints submitted to a federal commission over a decade, while allowing an ``undue administrative burden'' defense. This scenario sits on \textit{Garrett}'s borderline: public higher education is not a fundamental right and disability triggers only rational-basis review (favoring immunity under \textit{Garrett}), yet the record is targeted specifically at state universities rather than private actors and the remedy is narrower than Title I's workplace-wide mandate, leaving courts closely divided over congruence and proportionality.
\end{enumerate}

\section{Critique}

\textit{Board of Trustees of the University of Alabama v. Garrett} has provoked intense criticism across the ideological spectrum (\href{https://digitalcommons.law.umaryland.edu/cgi/viewcontent.cgi?article=3119&context=mlr}{Maryland Law Review}; \href{https://scholarship.kentlaw.iit.edu/cgi/viewcontent.cgi?article=3368&context=cklawreview}{Chicago-Kent Law Review}; \href{https://research.ualr.edu/bowen_lawreview/vol42/iss3/6/}{Bowen Law Review}).

\textbf{Progressive and Structural Constitutional Critique:} Progressive scholars argue that \textit{Garrett} exemplifies ``conservative judicial activism'' (\href{https://lawreview.colorado.edu/wp-content/uploads/2025/12/Young.pdf}{Colorado Law Review}) by converting Section 5 of the Fourteenth Amendment---an explicit constitutional grant of legislative power to Congress---into a judicially policed administrative-record test. Critics contend that demanding a judicial-style evidentiary record of unconstitutional state conduct retroactively penalizes Congress for failing to anticipate \textit{City of Boerne}'s 1997 standard when passing the ADA in 1990. Furthermore, scholars emphasize that excluding local-government discrimination ignores how deeply intertwined state and local employment systems are, and that leaving state employees with only federal DOJ enforcement or \textit{Ex parte Young} injunctions strips away the primary practical deterrent and make-whole remedy (back pay and damages) for unlawful termination.

\textbf{Originalist and Textualist Critique:} Surprisingly, \textit{Garrett} also draws sharp fire from rigorous textualists and originalists (\href{https://scholarship.law.uc.edu/cgi/viewcontent.cgi?article=1475&context=uclr}{University of Cincinnati Law Review}; \href{https://law.siu.edu/_common/documents/law-journal/articles-2016/fall-2016/6-sample-article-sm.pdf}{SIU Law Journal}). The Eleventh Amendment's plain text bars only federal suits against a State ``by Citizens of another State, or by Citizens or Subjects of any Foreign State''---not suits by a State's \textit{own} citizens like Patricia Garrett and Milton Ash. Textualists argue that \textit{Hans v. Louisiana} and \textit{Seminole Tribe} replaced precise constitutional text with an atextual, freestanding theory of state sovereign immunity, thereby undermining the textualist commitment to enforcing constitutional compromises as written. Conversely, some originalists defend \textit{Garrett}'s Section 5 holding on the ground that the Fourteenth Amendment's framers intended Section 5 to empower Congress to \textit{remedy} Section 1 violations, not to expand substantive constitutional rights beyond what Section 1 prohibits.

\textbf{Logical Weaknesses and Competing Values:} The decision exposes a deep doctrinal tension between judicial supremacy in constitutional interpretation (\textit{Marbury v. Madison} and \textit{City of Boerne}) and congressional primacy in enforcing Reconstruction's equality guarantees. By fusing \textit{Cleburne}'s ultra-deferential rational-basis test with \textit{Boerne}'s demanding congruence-and-proportionality inquiry, \textit{Garrett} created a near-insurmountable catch-22: because almost any state employment rule touching a non-suspect class can be hypothesized as ``rational'' in court, Congress can almost never prove a ``pattern'' of unconstitutional state conduct---effectively nullifying Section 5 abrogation for every group outside suspect or quasi-suspect classifications unless a fundamental right is implicated.

\section{Key Quotations}

\begin{quote}
``Thus, the result of \textit{Cleburne} is that States are not required by the Fourteenth Amendment to make special accommodations for the disabled, so long as their actions toward such individuals are rational. They could quite hardheadedly---and perhaps hardheartedly---hold to job-qualification requirements which do not make allowance for the disabled. If special accommodations for the disabled are to be required, they have to come from positive law and not through the Equal Protection Clause.''
\end{quote}
\noindent --- Chief Justice Rehnquist, Opinion of the Court (\href{https://supreme.justia.com/cases/federal/us/531/356/}{531 U.S. at 367--368}).

\begin{quote}
``Congress is the final authority as to desirable public policy, but in order to authorize private individuals to recover money damages against the States, there must be a pattern of discrimination by the States which violates the Fourteenth Amendment, and the remedy imposed by Congress must be congruent and proportional to the targeted violation. Those requirements are not met here, and to uphold the Act's application to the States would allow Congress to rewrite the Fourteenth Amendment law laid down by this Court in \textit{Cleburne}. Section 5 does not so broadly enlarge congressional authority.''
\end{quote}
\noindent --- Chief Justice Rehnquist, Opinion of the Court (\href{https://supreme.justia.com/cases/federal/us/531/356/}{531 U.S. at 374}).

\begin{quote}
``One of the undoubted achievements of statutes designed to assist those with impairments is that citizens have an incentive, flowing from a legal duty, to develop a better understanding, a more decent perspective, for accepting persons with impairments or disabilities into the larger society. The law works this way because the law can be a teacher. So I do not doubt that the Americans with Disabilities Act of 1990 will be a milestone on the path to a more decent, tolerant, progressive society.''
\end{quote}
\noindent --- Justice Kennedy, Concurring (\href{https://supreme.justia.com/cases/federal/us/531/356/}{531 U.S. at 375}).

\begin{quote}
``But a legislature is not a court of law. And Congress, unlike courts, must, and does, routinely draw general conclusions---for example, of likely motive or of likely relationship to legitimate need---from anecdotal and opinion-based evidence of this kind, particularly when the evidence lacks strong refutation.''
\end{quote}
\noindent --- Justice Breyer, Dissenting (\href{https://supreme.justia.com/cases/federal/us/531/356/}{531 U.S. at 379--380}).

\begin{quote}
``And it is difficult to understand why the Court, which applies `minimum ``rational-basis'' review' to statutes that burden persons with disabilities, \textit{ante}, at 366, subjects to far stricter scrutiny a statute that seeks to help those same individuals.''
\end{quote}
\noindent --- Justice Breyer, Dissenting (\href{https://supreme.justia.com/cases/federal/us/531/356/}{531 U.S. at 387--388}).


\end{document}
